% page 177 du fac-simile (image p-176 du scan)
% bloc 170.2 x 233.3 mm ; marge gauche 31.8 mm
\pgc{10.160mm}{0.000mm}{
\l{\cel{53}167}
\l{}
\l{}
\l{\sou{fiziko}. La cienco di qua la studiajo esas la propraji generala}
\l{\cel{3}di la korpi e la legi (+leyi) qui koncernas lia efiki extera,}
\l{\cel{3}sen egardar la kompozanti. - DEFIRS.}
\l{}
\l{\sou{fizognomonio}. La arto determinar la karaktero di persono segun}
\l{\cel{3}la traiti di lua vizajo. - DEFIRS.}
\l{}
\l{\sou{fiziokratio}. Doktrino ekonomikala qua konsideras ke tero (sulo)}
\l{\cel{3}esas la fonto unika di richeso. - DEFIRS.}
\l{}
\l{\sou{fiziologio}. La parto di la natur-historio qua traktas la roli}
\l{\cel{3}di la organi che la enti vivanta en lia stando normala. - DEFIRS}
\l{}
\l{\sou{fizionomio}. Expresuro che la vizajo. - DEFIRS.}
\l{}
\l{\sou{flado}. Tarto moleta, ek kremo, farino ed ovi. - D.}
\l{}
\l{\sou{flago}. Insigno (generale min granda kam standardo) qua indikas}
\l{\cel{3}la lojeyo, navo o loko, di ula chefo o di ula establisuro.}
\l{\cel{3}- DER.}
\l{}
\l{\sou{flagelo}. (\sou{histologio}).Filamento moviva qua (che ula protozoi, e}
\l{\cel{3}che la gameti maskula, ed anke che la celuli epiteliala) fun-}
\l{\cel{3}cionas kom organo lokomocala. - DEFIRS.}
\l{}
\l{\sou{flagrar}. (\sou{netrans.}) I. Kombustesar rapide kun flami granda e}
\l{\cel{3}subita. (\sou{anke metaf}.)- II. (\sou{pri delikto}). Eventar vidate da}
\l{\cel{3}la persono ipsa qua konstatas lu. - DEFIS.}
\l{}
\l{flajoleto. (\sou{muziko)}. Fluto bekoza, sis-tono, quan onu pluperfekt-}
\l{\cel{3}igas per adjuntar a lu klavi. - DEFS.}
\l{}
\l{\sou{flako}. Mikra marsheto de aquo, neprofunda. - F.}
\l{}
\l{\sou{flakono}. Mikra botelo, ek vitro, kristalo, porcelano, kun stopilo}
\l{\cel{3}generale vitra o metala. - DEFR.}
\l{}
\l{\sou{flamo}. Kombinuro de la oxo di aero kun la partikuli o la gasi qui}
\l{\cel{3}eskapas de la materii kombustata. - DEFIS.}
\l{}
\l{\sou{flamingo}. (\sou{zool.}) Ucelo de la klaso \textquotedbl{}steltieri\textquotedbl{}, di qua la suba}
\l{\cel{3}parto di la ali esas fairo-reda. - DEFIRS.}
\l{}
\l{flanar. (\sou{netrans}.) Irar sen skopo, lasante su distraktesar da ta}
\l{\cel{3}kozo, da ca kozo, por pasigar tempo. - DF.}
\l{}
\l{flanelo. Stofo dolca, lanugoza, ek lano pektita o kardita, di}
\l{\cel{3}qua la texuro esas kelke desdensa, simpla o krucumita.-DEFIRS.}
\l{}
\l{flanjo. (\sou{tekn.)}\cel{2}Parto diskatra,salianta, di tubi, stangi, e c.,}
\l{\cel{3}qua uzesas por ligar li per skrubi e bolti. - DEIR.}
}
